\documentclass[11pt,a4paper]{article}
\usepackage[utf8]{inputenc}
\usepackage[vietnamese]{babel}
\usepackage{amsmath,amssymb,amsthm}
\usepackage{enumitem}
\usepackage[margin=2.5cm]{geometry}

\newtheorem{theorem}{Định lý}
\newtheorem{assumption}{Giả thiết}
\newtheorem{definition}{Định nghĩa}
\newtheorem{lemma}{Bổ đề}
\newtheorem{remark}{Nhận xét}
\newtheorem{corollary}{Hệ quả}
\newtheorem{proposition}{Mệnh đề}

\newcommand{\erank}{\mathrm{erank}}
\newcommand{\rank}{\mathrm{rank}}
\newcommand{\RR}{\mathbb{R}}
\newcommand{\EE}{\mathbb{E}}
\newcommand{\cL}{\mathcal{L}}
\newcommand{\cS}{\mathcal{S}}
\newcommand{\Tr}{\mathrm{Tr}}
\newcommand{\diag}{\mathrm{diag}}

\title{Chứng minh chi tiết Theorem 2: \\ Bảo toàn hạng hiệu dụng dưới cập nhật có tiền điều kiện}
\author{}
\date{}

\begin{document}
\maketitle

%=======================================================================
\section{Ký hiệu và thiết lập}
%=======================================================================

\paragraph{Ma trận trọng số và gradient.}
Xét một lớp có ma trận trọng số $W \in \RR^{m \times n}$ và gradient tại một bước là
\[
G = \nabla_W \cL \in \RR^{m \times n}.
\]
Một bước cập nhật có tiền điều kiện có dạng
\[
W_{t+1} = W_t - \eta P_t^{-1}G_t,
\]
trong đó $P_t \succ 0$ là preconditioner và $\eta > 0$ là learning rate.

\paragraph{Phân tích SVD.}
Với một ma trận $X \in \RR^{m \times n}$, viết SVD rút gọn
\[
X = U_X \Sigma_X V_X^\top,
\qquad
\Sigma_X = \diag(\sigma_1(X), \ldots, \sigma_r(X)),
\qquad
r = \min(m,n),
\]
trong đó
\[
\sigma_1(X) \geq \sigma_2(X) \geq \cdots \geq \sigma_r(X) \geq 0.
\]
Ta ký hiệu condition number phổ của $X$ bởi
\[
\kappa(X) = \frac{\sigma_1(X)}{\sigma_r(X)}
\]
khi $X$ có hạng đầy đủ trên không gian đang xét.

\begin{definition}[Hạng hiệu dụng dựa trên năng lượng phổ]
\label{def:erank}
Cho $X \in \RR^{m \times n}$ có singular values $\sigma_1,\dots,\sigma_r$, đặt
\[
p_i(X) = \frac{\sigma_i(X)^2}{\sum_{j=1}^r \sigma_j(X)^2} = \frac{\sigma_i(X)^2}{\|X\|_F^2}.
\]
Hạng hiệu dụng của $X$ được định nghĩa bởi
\[
\erank(X) = \exp\!\left(-\sum_{i=1}^{r} p_i(X) \log p_i(X)\right).
\]
Đây là hàm mũ của entropy Shannon của phân phối năng lượng phổ.
\end{definition}

\paragraph{Mục tiêu của file này.}
Theorem 2 trong \texttt{paper\_merged.tex} có ba ý. Sau khi kiểm tra proof sketch trong appendix, phần một-bước về $P^{-1}G$ có thể chứng minh được nếu thêm giả thiết căn chỉnh phổ; phần dài hạn cho $W_T$ cần giả thiết động lực học bổ sung. File này viết lại phiên bản chứng minh được một cách chặt chẽ.

%=======================================================================
\section{Các giả thiết}
%=======================================================================

\begin{assumption}[Căn chỉnh phổ cục bộ giữa $G$ và $P$]
\label{ass:alignment}
Xét một bước cập nhật cố định. Tồn tại các ma trận trực giao $U \in \RR^{m \times r}$, $V \in \RR^{n \times r}$ sao cho
\[
G = U\diag(s_1,\ldots,s_r)V^\top,
\qquad
s_1 \ge s_2 \ge \cdots \ge s_r > 0,
\]
và preconditioner $P$ tác động lên cùng basis trái đó:
\[
P U = U\diag(\lambda_1,\ldots,\lambda_r),
\qquad
\lambda_1 \ge \lambda_2 \ge \cdots \ge \lambda_r > 0.
\]
Nói cách khác, các hướng phổ trội của gradient cũng là các hướng curvature lớn mà $P^{-1}$ sẽ làm giảm mạnh hơn.
\end{assumption}

\begin{assumption}[Không đảo thứ tự sau preconditioning]
\label{ass:no_reorder}
Sau khi chia theo các hệ số $\lambda_i$, thứ tự giảm của singular values vẫn được giữ:
\[
\frac{s_i}{\lambda_i} \ge \frac{s_j}{\lambda_j}
\qquad \text{với mọi } i<j.
\]
Tương đương,
\[
\frac{s_i^2}{s_j^2} \ge \frac{\lambda_i^2}{\lambda_j^2}
\qquad \text{với mọi } i<j.
\]
\end{assumption}

\begin{assumption}[Kích thích cân bằng sau preconditioning]
\label{ass:balanced_excitation}
Tồn tại một cơ sở trực chuẩn cố định $\{E_1,\dots,E_r\}$ trong không gian hoạt động của các cập nhật đã precondition sao cho
\[
P_t^{-1}G_t = \sum_{i=1}^r \alpha_{t,i}E_i.
\]
Đặt
\[
A_i(T) = \sum_{t=0}^{T-1} \alpha_{t,i}.
\]
Ta giả sử tồn tại hằng số $C \ge 1$ sao cho với $T$ đủ lớn,
\[
\frac{1}{Cr} \le \frac{A_i(T)^2}{\sum_{j=1}^r A_j(T)^2} \le \frac{C}{r}
\qquad \text{cho mọi } i=1,\dots,r,
\]
và phần cập nhật tích lũy trội hơn khởi tạo ban đầu:
\[
\|W_0\|_F = o\!\left(\left\|\sum_{t=0}^{T-1} P_t^{-1}G_t\right\|_F\right).
\]
\end{assumption}

\begin{assumption}[Learning rate cố định]
\label{ass:lr}
Learning rate $\eta > 0$ là hằng số.
\end{assumption}

\begin{remark}
Assumption Gradient Concentration trong paper là chưa đủ để suy ra trực tiếp các bất đẳng thức entropy trong Theorem 2. Điểm còn thiếu là mối liên hệ hình học giữa basis của $G$ và basis của $P$. Hai Giả thiết~\ref{ass:alignment} và~\ref{ass:no_reorder} chính là phần cần bổ sung để chứng minh một-bước trở nên chặt.
\end{remark}

%=======================================================================
\section{Phát biểu định lý}
%=======================================================================

\begin{theorem}[Phiên bản chặt của phần một-bước trong Theorem 2]
\label{thm:rank_preservation_local}
Dưới các Giả thiết~\ref{ass:alignment}, \ref{ass:no_reorder} và \ref{ass:lr}, với cập nhật
\[
W_{t+1} = W_t - \eta P_t^{-1}G_t,
\]
ta có:
\begin{enumerate}[label=(\roman*)]
    \item \textbf{Co condition number của update:}
    \[
    \kappa(P^{-1}G) = \frac{\kappa(G)}{\kappa(P)}.
    \]

    \item \textbf{Hạng hiệu dụng của update không giảm:}
    \[
    \erank(P^{-1}G) \ge \erank(G).
    \]
\end{enumerate}
\end{theorem}

\begin{corollary}[Cận dưới định lượng theo condition number]
\label{cor:entropy_bound}
Dưới Giả thiết~\ref{ass:alignment}, nếu đặt
\[
K = \frac{\kappa(G)}{\kappa(P)},
\qquad
a = \frac{K^2}{K^2+r-1},
\qquad
b = \frac{1}{K^2+r-1},
\]
thì
\[
\erank(P^{-1}G) \ge \exp\!\bigl(-a\log a -(r-1)b\log b\bigr).
\]
\end{corollary}

\begin{proposition}[Phiên bản dài hạn cho $W_T$]
\label{prop:long_run}
Dưới thêm Giả thiết~\ref{ass:balanced_excitation}, tồn tại hằng số $c = c(C) > 0$ sao cho với $T$ đủ lớn,
\[
\erank(W_T) \ge c r.
\]
Cụ thể hơn, có thể lấy $c = 1/C$ ở mức cận dưới thô.
\end{proposition}

%=======================================================================
\section{Chứng minh chi tiết}
%=======================================================================

\subsection{Bước 1: Viết rõ singular values của update đã precondition}

Từ Giả thiết~\ref{ass:alignment}, ta có
\[
G = U\diag(s_1,\ldots,s_r)V^\top
\quad \text{và} \quad
P^{-1}U = U\diag(\lambda_1^{-1},\ldots,\lambda_r^{-1}).
\]
Do đó
\[
P^{-1}G
= U\diag(\lambda_1^{-1},\ldots,\lambda_r^{-1})U^\top U\diag(s_1,\ldots,s_r)V^\top
= U\diag(t_1,\ldots,t_r)V^\top,
\]
trong đó
\[
t_i = \frac{s_i}{\lambda_i}, \qquad i=1,\ldots,r.
\]
Vì vậy, singular values của update đã precondition chính là singular values của gradient sau khi được rescale từng hướng bởi $1/\lambda_i$.

\subsection{Bước 2: Chứng minh phần (i) của định lý}

Ta có
\[
\kappa(P^{-1}G) = \frac{t_1}{t_r}
= \frac{s_1/\lambda_1}{s_r/\lambda_r}
= \frac{s_1}{s_r}\cdot\frac{\lambda_r}{\lambda_1}
= \frac{\kappa(G)}{\kappa(P)}.
\]
Vậy phần (i) được chứng minh. \qed

\subsection{Bước 3: Chuẩn hóa năng lượng phổ trước và sau preconditioning}

Đặt
\[
p_i = \frac{s_i^2}{\sum_{j=1}^r s_j^2},
\qquad
\tilde p_i = \frac{t_i^2}{\sum_{j=1}^r t_j^2}
= \frac{s_i^2/\lambda_i^2}{\sum_{j=1}^r s_j^2/\lambda_j^2}.
\]
Khi đó
\[
\erank(G) = \exp\!\left(-\sum_{i=1}^r p_i\log p_i\right),
\qquad
\erank(P^{-1}G) = \exp\!\left(-\sum_{i=1}^r \tilde p_i\log \tilde p_i\right).
\]

Mục tiêu là chứng minh phân phối năng lượng phổ mới $\tilde p$ ``đều hơn'' phân phối cũ $p$.

\subsection{Bước 4: Một bổ đề reweighting đơn điệu}

\begin{lemma}[Reweighting bởi dãy tăng làm đều phân phối]
\label{lem:reweight_majorization}
Cho $p_1 \ge p_2 \ge \cdots \ge p_r \ge 0$ với $\sum_i p_i = 1$. Cho thêm một dãy trọng số không giảm
\[
b_1 \le b_2 \le \cdots \le b_r, \qquad b_i > 0,
\]
và định nghĩa phân phối mới
\[
q_i = \frac{p_i b_i}{\sum_{j=1}^r p_j b_j}.
\]
Khi đó $q \prec p$, tức là $q$ bị majorize bởi $p$.
\end{lemma}

\begin{proof}
Ta cần kiểm tra: với mọi $k = 1, \ldots, r-1$,
\[
\sum_{i=1}^k q_i \leq \sum_{i=1}^k p_i.
\]
Vì $p_i$ giảm và theo Giả thiết~\ref{ass:no_reorder} thì $q_i$ cũng giảm, ta không cần sắp xếp lại; chỉ cần kiểm tra tổng đầu trực tiếp:
\[
\sum_{i=1}^k q_i = \frac{\sum_{i=1}^k p_i b_i}{\sum_{j=1}^r p_j b_j}.
\]
Ta muốn chỉ ra điều này $\leq \sum_{i=1}^k p_i$, hay tương đương (nhân chéo, mẫu dương):
\[
\biggl(\sum_{i=1}^k p_i b_i\biggr)\biggl(\sum_{j=k+1}^r p_j\biggr)
\;\leq\;
\biggl(\sum_{i=1}^k p_i\biggr)\biggl(\sum_{j=k+1}^r p_j b_j\biggr).
\tag{$\star$}
\]
Chia hai vế cho $(\sum_{i=1}^k p_i)(\sum_{j=k+1}^r p_j) > 0$, bất đẳng thức $(\star)$ trở thành:
\[
\underbrace{\frac{\sum_{i=1}^k p_i b_i}{\sum_{i=1}^k p_i}}_{\text{trung bình có trọng số của } b \text{ trên nhóm đầu}}
\;\leq\;
\underbrace{\frac{\sum_{j=k+1}^r p_j b_j}{\sum_{j=k+1}^r p_j}}_{\text{trung bình có trọng số của } b \text{ trên nhóm đuôi}}.
\]
Vì $b$ không giảm ($b_1 \leq \cdots \leq b_r$):
\begin{itemize}[nosep]
    \item Nhóm đầu: mọi $b_i \leq b_k$ khi $i \leq k$ $\;\Longrightarrow\;$ trung bình $\leq b_k$.
    \item Nhóm đuôi: mọi $b_j \geq b_{k+1} \geq b_k$ khi $j > k$ $\;\Longrightarrow\;$ trung bình $\geq b_k$.
\end{itemize}
Do đó vế trái $\leq b_k \leq$ vế phải, bất đẳng thức $(\star)$ đúng, suy ra $\sum_{i=1}^k q_i \leq \sum_{i=1}^k p_i$ với mọi $k$. Kết hợp với $\sum_{i=1}^r q_i = \sum_{i=1}^r p_i = 1$, ta được $q \prec p$.
\end{proof}

\subsection{Bước 5: Chứng minh phần (ii) của định lý}

Áp dụng Bổ đề~\ref{lem:reweight_majorization} với
\[
b_i = \lambda_i^{-2}.
\]
Do $\lambda_1 \ge \cdots \ge \lambda_r > 0$, ta có
\[
\lambda_1^{-2} \le \lambda_2^{-2} \le \cdots \le \lambda_r^{-2}.
\]
Hơn nữa,
\[
\tilde p_i = \frac{p_i \lambda_i^{-2}}{\sum_{j=1}^r p_j \lambda_j^{-2}}.
\]
Vậy theo Bổ đề~\ref{lem:reweight_majorization},
\[
\tilde p \prec p.
\]
Vì entropy Shannon là hàm Schur-concave, ta suy ra
\[
H(\tilde p) \ge H(p).
\]
Lấy hàm mũ hai vế,
\[
\erank(P^{-1}G) = e^{H(\tilde p)} \ge e^{H(p)} = \erank(G).
\]
Phần (ii) được chứng minh. \qed

\subsection{Bước 6: Chứng minh hệ quả định lượng theo condition number}

Trong lớp các ma trận hạng $r$ có condition number không vượt quá $K$, cấu hình có entropy nhỏ nhất là cấu hình cực trị
\[
(K,1,\ldots,1).
\]
Sau khi chuẩn hóa theo bình phương singular values, ta được phân phối
\[
(a,b,\ldots,b),
\qquad
a = \frac{K^2}{K^2+r-1},
\qquad
b = \frac{1}{K^2+r-1}.
\]
Do đó mọi ma trận $X$ thỏa $\kappa(X) \le K$ đều phải có
\[
\erank(X) \ge \exp\!\bigl(-a\log a -(r-1)b\log b\bigr).
\]
Áp dụng với $X = P^{-1}G$ và dùng phần (i) cho $K = \kappa(G)/\kappa(P)$, ta thu được Hệ quả~\ref{cor:entropy_bound}. \qed

\subsection{Bước 7: Chứng minh mệnh đề dài hạn cho $W_T$}

Xét phần cập nhật tích lũy
\[
\Delta_T = \sum_{t=0}^{T-1} P_t^{-1}G_t = \sum_{i=1}^r A_i(T)E_i.
\]
Theo Giả thiết~\ref{ass:balanced_excitation}, với $T$ đủ lớn,
\[
\frac{1}{Cr} \le \frac{A_i(T)^2}{\sum_{j=1}^r A_j(T)^2} \le \frac{C}{r}.
\]
Điều đó có nghĩa là năng lượng của $\Delta_T$ trên cơ sở $\{E_i\}$ được phân bố gần đều ở cấp độ hằng số, không bị dồn vào một số hữu hạn ít hướng.

Đặt
\[
q_i(T) = \frac{A_i(T)^2}{\sum_{j=1}^r A_j(T)^2}.
\]
Khi đó
\[
q_i(T) \le \frac{C}{r}
\qquad \text{cho mọi } i.
\]
Suy ra entropy của phân phối $q(T)$ có cận dưới thô:
\[
H(q(T)) = -\sum_{i=1}^r q_i(T)\log q_i(T)
\ge -\sum_{i=1}^r q_i(T)\log\!\left(\frac{C}{r}\right)
= \log\!\left(\frac{r}{C}\right).
\]
Do đó
\[
\erank(\Delta_T) \ge \frac{r}{C}.
\]

Bây giờ dùng giả thiết
\[
\|W_0\|_F = o(\|\Delta_T\|_F),
\]
ta có
\[
W_T = W_0 - \eta \Delta_T
\]
và vì thành phần $-\eta\Delta_T$ trội hơn $W_0$ khi $T$ lớn, phân phối năng lượng phổ của $W_T$ tiệm cận phân phối năng lượng phổ của $\Delta_T$ ở mức hằng số. Vì thế tồn tại hằng số $c = c(C) > 0$ sao cho với $T$ đủ lớn,
\[
\erank(W_T) \ge c r.
\]
Một cận dưới thô có thể lấy là $c = 1/C$ sau khi hấp thụ sai lệch hữu hạn do $W_0$ vào hằng số. Mệnh đề được chứng minh. \qed

%=======================================================================
\section{Kết luận kiểm chứng}
%=======================================================================

\begin{remark}
Từ chứng minh trên, có ba kết luận quan trọng.
\begin{enumerate}[label=(\roman*)]
    \item Phần một-bước của Theorem 2 có thể chứng minh được, nhưng cần thêm giả thiết căn chỉnh phổ giữa $G$ và $P$.
    \item Điều kiện chỉ dựa vào $\kappa(P)$ trong bản sketch của paper là chưa đủ nếu không nói rõ cấu trúc hình học của basis.
    \item Phần dài hạn cho $W_T$ không thể suy ra chỉ từ câu ``mỗi update có hạng hiệu dụng cao hơn''; nó cần giả thiết động lực học bổ sung như Giả thiết~\ref{ass:balanced_excitation}.
\end{enumerate}
Vì vậy, nếu muốn giữ paper chặt về mặt toán học, nên tách Theorem 2 thành: một kết quả một-bước cho $P^{-1}G$ và một mệnh đề dài hạn riêng cho $W_T$.
\end{remark}

\end{document}
